\documentclass[10pt,a4paper]{amsart}
\usepackage[margin=19mm]{geometry}
\usepackage{amsmath,amssymb,mathtools}
\usepackage[scr=rsfs,cal=euler]{mathalfa}
\usepackage{xfrac}
\usepackage[unicode,hidelinks,bookmarks=false]{hyperref}
\newcommand{\ie}{i.e.\ }
\newcommand{\cf}{cf.\ }
\newcommand{\resp}{respectively\ }
\newcommand{\Subsection}{\subsection}
\providecommand{\nobreakdash}{\nobreak\hbox{-}}
\setlength{\emergencystretch}{2em}
\multlinegap0pt
\usepackage{bm}
\usepackage{bbm}
\usepackage{dsfont}
\usepackage{xspace}
\usepackage{cancel}
\usepackage{mathtools}
\usepackage{amsthm}
\makeatletter
\@ifundefined{theorem}{\newtheorem{theorem}{Theorem}}{}
\@ifundefined{definition}{\newtheorem{definition}[theorem]{Definition}}{}
\@ifundefined{lemma}{\newtheorem{lemma}[theorem]{Lemma}}{}
\makeatother
\def\Bar{\overline}
\newcommand{\R}{\mathbb{R}}
\title[Solving Regularized Multifacility Location Problems with Unk]{Solving Regularized Multifacility Location Problems with Unknown Number of Centers via Difference-of-Convex Optimization}
\author{W. Geremew, V. S. T. Long, N. M. Nam, A. Solano-Herrera}
\date{Source: arXiv:2601.16576v1 (CC BY 4.0)}
\begin{document}
\maketitle

\noindent\emph{Source and licence.} Solving Regularized Multifacility Location Problems with Unknown Number of Centers via Difference-of-Convex Optimization. Version arXiv:2601.16576v1 (CC BY 4.0), distributed under the Creative Commons Attribution 4.0 International licence, as recorded in the arXiv licence metadata for this exact version and shown on the abstract page https://arxiv.org/abs/2601.16576v1. Full rights notice: licenses/arxiv-2601.16576-CC-BY-4.0.txt.\par\smallskip

\noindent\textit{Editorial scope.} Counted unit: the Theorem whose source label is \texttt{thm:local} in the article (source0.tex lines 720--727, source label \texttt{thm:local}), delivered together with the article's own complete original proof of it (source0.tex lines 729--828, one \texttt{proof} environment of 100 lines). No mathematics is written, repaired or re-derived: statement and proof are the article's own text line for line. Required context retained uncounted: maxminn (lines 77--79); deflocal (lines 622--634); lem:Li-stable (lines 661--673). Every definition, display and label of those lines is retained unchanged. Omitted from this package and not redistributed: 1--76, 80--621, 635--660, 674--719, 728--728, 829--2561. Nothing inside a retained excerpt is omitted or altered: every retained body is delivered exactly as the article's lines are written, so no omission receipt of any kind accompanies this package. Citations: the retained excerpts carry no citation command at all, so no bibliography entry of the article is transcribed and no reference is redistributed. Reference adaptation: none was needed and none was made; every reference inside the delivered unit resolves inside it, so no unresolved reference or citation is produced. Printed slips kept literal: no printed word, symbol, hypothesis or number of the counted statement or of its proof was altered. Layout and class: the article's own class is \texttt{article} with packages including geometry, times, verbatim, amsmath, amssymb, amsthm, mathtools, bbm, bm, xcolor, graphicx, subcaption, algorithm, algpseudocode; the wrapper is the lane's pinned \texttt{amsart} editorial wrapper at 10pt on A4, which restates the theorem-like environments the retained text needs and adds the article's own macro definitions that the retained text needs (\texttt{\textbackslash Bar}, \texttt{\textbackslash R}), with extra wrapper packages (bm, bbm, dsfont, xspace, cancel, mathtools). The delivered numbering is the unit's own. Assets: no retained span contains a figure environment, a graphics-inclusion command, a PDF-page inclusion, a table or a TikZ picture; no image file is copied into this package, the package declares no asset dependency and no image file is redistributed with it. Licence: version arXiv:2601.16576v1 of this article is distributed under the Creative Commons Attribution 4.0 International licence, as recorded in the arXiv licence metadata for this exact version and shown on the abstract page https://arxiv.org/abs/2601.16576v1. A full rights notice is delivered at licenses/arxiv-2601.16576-CC-BY-4.0.txt. Characters: every retained excerpt is pure ASCII, so the delivered engine, XeLaTeX with its default Latin Modern text font, reports no missing character.
\begin{equation}\label{maxminn}
	\min_{X\in\mathbb{R}^{k\times d}} \Biggl\{ f(X) = \sum_{i=1}^{n} \min_{1 \le \ell \le k} \rho(x_\ell - a_i) + \frac{\lambda n}{2} \sum_{1 \le s < t \le k} \|x_s - x_t\|^2 \Biggr\},
\end{equation}

\begin{definition}\label{deflocal}
An element $\Bar X=(\bar x_1,\dots,\bar x_k)\in\R^{k\times d}$ is called a
\emph{local optimal solution} of problem~\eqref{maxminn} if there exists
$\varepsilon>0$ such that
\[
f(\Bar X)\;\le\;f(X)
\quad
\text{for all }X=(x_1,\dots,x_k)\in\R^{k\times d}
\text{ with }\|x_\ell-\bar x_\ell\|<\varepsilon
\text{ for all }\ell.
\]
The set of all local optimal solutions will be denoted by $S_{loc}$.
\end{definition}

\begin{lemma}\label{lem:Li-stable}
Let $\Bar X=(\bar x_1,\dots,\bar x_k)\in\R^{k\times d}$ be such that
$L_i(\Bar X)$ is a singleton for every $i=1,\dots,n$. Then there exists
$\varepsilon>0$ such that for every
$X=(x_1,\dots,x_k)\in\R^{k\times d}$ satisfying
$\|x_\ell-\bar x_\ell\|<\varepsilon$ for all $\ell=1,\dots,k$, we have:
\begin{itemize}
    \item [\rm{(a)}] $L_i(X) = L_i(\Bar X)\quad\text{for all } i=1,\dots,n.$
    \item [\rm{(b)}] 
$I_\ell(X) = I_\ell(\Bar X)\quad\text{and}\quad A_\ell(X)=A_\ell(\Bar X)
\quad\text{for all } \ell=1,\dots,k.$
\end{itemize}
\end{lemma}

\begin{theorem}\label{thm:local}
Let $\Bar X=(\bar x_1,\dots,\bar x_k)\in\mathbb{R}^{k\times d}$ be such that $L_i(\Bar X)$ is a singleton for every $i=1,\dots,n$. If $\Bar X$ is a local optimal solution of problem~\eqref{maxminn}, then for each $\ell\in\{1,\dots,k\}$, the point $\bar x_\ell$ is a global minimizer of the problem
\begin{equation}\label{eq:phi-ell}
\min_{x\in\mathbb{R}^d} \Biggl\{
\phi_\ell(x) = \sum_{i\in I_\ell(\Bar X)}\rho(x-a_i) + \frac{\lambda n}{2} \sum_{j=1,\; j \neq \ell}^k \|x-\bar x_j\|^2
\Biggr\}.
\end{equation}
\end{theorem}

\begin{proof}
By Lemma~\ref{lem:Li-stable}, the assumption that $L_i(\Bar X)$ is a singleton
for every $i$ implies the following stability property, i.e., there exists
$\varepsilon>0$ such that, for all $X=(x_1,\dots,x_k)\in\R^{k\times d}$ with
\begin{equation}\label{eq4}
\|x_\ell-\bar x_\ell\|<\varepsilon\qquad\text{for all }\ell=1,\dots,k,
\end{equation}
we have
\begin{equation}\label{eq5}
L_i(X)=L_i(\Bar X)\qquad\text{for all }i=1,\dots,n.
\end{equation}
In particular, in the neighbourhood defined by~\eqref{eq4}, we have
\begin{equation}\label{key15}
    I_\ell(X)=I_\ell(\Bar X)
\quad\text{and}\quad
A_\ell(X)=A_\ell(\Bar X)
\qquad\text{for all }\ell=1,\dots,k.
\end{equation}
Suppose now that $\Bar X$ is a local optimal solution of~\eqref{maxminn}. 
By Definition~\ref{deflocal}, there exists $\delta>0$ such that
\[
f(\Bar X)\;\le\;f(X)
\quad
\text{for all }X=(x_1,\dots,x_k)\in\R^{k\times d}
\text{ with }\|x_\ell-\bar x_\ell\|<\delta
\text{ for all }\ell.
\]
Let $\varepsilon>0$ be as in~\eqref{eq4}, and set $
\gamma = \min\{\varepsilon,\delta\}.$
 For a fixed index $\ell$, take any $x\in\mathbb{R}^d$ such that $\|x-\bar x_\ell\| < \gamma$ and define 
 \[
\widehat x_\ell = x
\quad\text{and}\quad
\widehat x_r = \bar x_r \ \text{ for } r\neq\ell.
\]
Then $\widehat X$ satisfies \eqref{eq4}, and by local optimality we have
\begin{equation}\label{eq6}
f(\Bar X)\le f(\widehat X).
\end{equation}
To compare these two values, write
\begin{equation}\label{fDC}
f(X)=D(X)+C(X),    
\end{equation}
where
\[
D(X)=\sum_{i=1}^n\min_{1\le \ell\le k}\rho(x_\ell-a_i)
\quad\text{and}\quad
C(X)=\frac{\lambda n}{2}\sum_{1\le s<t\le k}\|x_s-x_t\|^2.
\]
 By~\eqref{eq5} and \eqref{key15}, we have
\begin{equation*}\label{eq:D-bar}
D(\Bar X)
=
\sum_{i\in I_\ell(\Bar X)}\rho(\bar x_\ell-a_i)
+
\sum_{i\notin I_\ell(\Bar X)}\min_{1\le r\le k}\rho(\bar x_r-a_i),
\end{equation*}
and
\begin{equation*}\label{eq:D-hat}
\begin{aligned}
     D(\widehat X)
     &=
\sum_{i\in I_\ell(\widehat X)}\rho(x-a_i)
+
\sum_{i\notin I_\ell(\widehat X)}\min_{1\le r\le k}\rho(\bar x_r-a_i) \\
     &=
\sum_{i\in I_\ell(\Bar X)}\rho(x-a_i)
+
\sum_{i\notin I_\ell(\Bar X)}\min_{1\le r\le k}\rho(\bar x_r-a_i),
\end{aligned}
\end{equation*}
which implies that
\begin{equation}\label{eq:D-diff}
D(\widehat X) - D(\Bar X) = \sum_{i\in I_\ell(\Bar X)} \rho(x-a_i) - \sum_{i\in I_\ell(\Bar X)} \rho(\bar x_\ell-a_i).
\end{equation}

For the fusion term, we can group as follows:
\begin{equation}\label{eq:C-bar-new}
C(\Bar{X}) = \frac{\lambda n}{2} \left( \sum_{j=1,\; j \neq \ell}^k \|\bar{x}_\ell - \bar{x}_j\|^2 + \sum_{1 \le s < t \le k,\; s,t \neq \ell} \|\bar{x}_s - \bar{x}_t\|^2 \right).
\end{equation}
Similarly for $\widehat{X}$, we have
\begin{equation}\label{eq:C-hat-new}
C(\widehat{X}) = \frac{\lambda n}{2} \left( \sum_{j=1,\; j \neq \ell}^k \|x - \bar{x}_j\|^2 + \sum_{1 \le s < t \le k,\; s,t \neq \ell} \|\bar{x}_s - \bar{x}_t\|^2 \right).
\end{equation}
Subtracting \eqref{eq:C-bar-new} from \eqref{eq:C-hat-new} yields
\begin{equation*}\label{eq:C-diff}
C(\widehat{X}) - C(\Bar{X}) = \frac{\lambda n}{2} \sum_{j=1,\; j \neq \ell}^k \left( \|x - \bar{x}_j\|^2 - \|\bar{x}_\ell - \bar{x}_j\|^2 \right).
\end{equation*}
Combining this with \eqref{fDC} and
\eqref{eq:D-diff}, we obtain
\begin{equation}
f(\widehat X) - f(\Bar X) = \phi_\ell(x) - \phi_\ell(\bar x_\ell),
\end{equation}
where $\phi_\ell$ is defined in \eqref{eq:phi-ell}. It follows from \eqref{eq6} that
$\phi_\ell(\bar x_\ell) \le \phi_\ell(x)$ for all $x$ in a neighborhood of $\bar x_\ell$. 
Since $\phi_\ell$ is convex,
every local minimizer is a global minimizer. Hence
$\bar x_\ell$ is a global solution of~\eqref{eq:phi-ell} for every $\ell$,
which proves the theorem.
\end{proof}

\end{document}
